% Three separate decision/validation processes. The only cross-process data
% edge is candidate evidence generated by a fixed EviComRL policy.
\begin{tikzpicture}[
  x=1mm,y=1mm,>=Stealth,
  font=\fontsize{9bp}{11.5bp}\selectfont,
  rfbox/.style={draw,line width=.45pt,fill=white,align=center,
    text width=36mm,minimum height=14mm,inner sep=1mm},
  rflabel/.style={align=center,text width=10mm,inner sep=0pt},
  rfhead/.style={align=center,font=\fontsize{10bp}{13bp}\selectfont},
  rfvalidation/.style={draw,line width=.35pt,fill=black!3,align=center,
    text width=128mm,minimum height=4.2mm,inner sep=.4mm},
  rfarrow/.style={->,line width=.55pt},
  rfreturn/.style={->,line width=.4pt},
  rfedgeword/.style={fill=white,inner sep=.25mm}
]
\path[use as bounding box] (0,0) rectangle (148,-90);
\node[rfhead] at (33.5,-2.5) {状态输入};
\node[rfhead] at (79.5,-2.5) {算法与动作};
\node[rfhead] at (125.5,-2.5) {结果与状态更新};
\foreach \yt/\yb in {-7/-32,-34/-59,-65/-90}
{
  \draw[line width=.45pt] (0,\yt) rectangle (148,\yb);
  \draw[line width=.35pt] (12,\yt) -- (12,\yb);
}
\node[rflabel] at (6,-19.5) {研究一\\网络\\传输};
\node[rflabel] at (6,-46.5) {研究二\\内容\\选择};
\node[rflabel] at (6,-77.5) {研究三\\语义\\交付};

\node[rfbox] (rfain) at (33.5,-15)
  {数据队列与剩余电量\\邻居链路及相对位置\\虚拟能量队列};
\node[rfbox] (rfadecision) at (79.5,-15)
  {\textbf{LEADR}\\时隙优化与路由评分\\下一跳与发送功率};
\node[rfbox] (rfaout) at (125.5,-15)
  {链路服务与成功转发\\数据队列与电量更新\\HELLO 邻居表维护};
\draw[rfarrow] (rfain.east) -- (rfadecision.west);
\draw[rfarrow] (rfadecision.east) -- (rfaout.west);
\draw[rfreturn] (rfaout.south) -- (125.5,-24.5) -- (33.5,-24.5) -- (rfain.south);
\node[rfedgeword] at (79.5,-24.5) {下一时隙状态};
\node[rfvalidation] at (79.5,-29.2)
  {网络仿真：节点数、通信范围、速度；吞吐、丢包、能量及寿命};

\node[rfbox] (rfbin) at (33.5,-42)
  {局部证据与链路代价\\中心信念及已选历史\\通信预算与动作掩码};
\node[rfbox] (rfbdecision) at (79.5,-42)
  {\textbf{EviComRL}\\图结构历史感知 PPO\\来源、粒度或停止};
\node[rfbox] (rfbout) at (125.5,-42)
  {已选语义候选\\中心信念与历史更新\\通信预算扣减};
\draw[rfarrow] (rfbin.east) -- (rfbdecision.west);
\draw[rfarrow] (rfbdecision.east) -- (rfbout.west);
\draw[rfreturn] (rfbout.south) -- (125.5,-51.5) -- (33.5,-51.5) -- (rfbin.south);
\node[rfedgeword] at (79.5,-51.5) {下一消息步骤状态};
\node[rfvalidation] at (79.5,-56.2)
  {语义轨迹：固定证据与代价；分类质量、有效信念质量与通信成本};

\node[rfbox] (rfcin) at (33.5,-73)
  {候选证据与链路预测\\发送记录和累计确认\\时限、电量与能量队列};
\node[rfbox] (rfcdecision) at (79.5,-73)
  {\textbf{C-BFC}\\效用、信念与两步前瞻\\新发、重传、等待、停止};
\node[rfbox] (rfcout) at (125.5,-73)
  {发送与有效确认\\接收信念及资源更新\\终止分类与通信代价};
\draw[rfarrow] (rfcin.east) -- (rfcdecision.west);
\draw[rfarrow] (rfcdecision.east) -- (rfcout.west);
\draw[rfreturn] (rfcout.south) -- (125.5,-82.5) -- (33.5,-82.5) -- (rfcin.south);
\node[rfedgeword] at (79.5,-82.5) {确认与资源状态};
\node[rfvalidation] at (79.5,-87.2)
  {事件回放：信道、反馈与时限；按时正确完成率、单位正确任务能耗、重复发送与时延};

% Content hand-off, with no online joint training or RQ1 data link implied.
\draw[rfarrow] (rfbout.east) -- (146,-42) -- (146,-61.7) -- (33.5,-61.7) -- (rfcin.north);
\node[rfedgeword] at (88,-61.7) {固定策略候选：合并去重至多四条};
\end{tikzpicture}
